\documentclass[11pt,letterpaper]{article}
\usepackage[margin=0.75in]{geometry}
\usepackage{enumitem}
\usepackage{titlesec}
\usepackage{hyperref}
\usepackage{xcolor}
\usepackage{fancyhdr}
\usepackage{lastpage}

\definecolor{linkblue}{HTML}{1a5276}

\hypersetup{
    colorlinks=true,
    linkcolor=linkblue,
    urlcolor=linkblue,
}

\pagestyle{fancy}
\fancyhf{}
\renewcommand{\headrulewidth}{0pt}
\rfoot{\thepage/\pageref{LastPage}}

\titleformat{\section}{\large\bfseries}{}{0em}{}[\titlerule]
\titlespacing{\section}{0pt}{12pt}{6pt}

\titleformat{\subsection}{\bfseries\color{linkblue}}{}{0em}{}
\titlespacing{\subsection}{0pt}{12pt}{4pt}

\setlist[itemize]{leftmargin=1.5em, itemsep=2pt, parsep=0pt, topsep=2pt}

\newcommand{\entrytitle}[2]{%
  \noindent\textbf{#1}\hfill #2\\
}

\newcommand{\entrysubtitle}[1]{%
  \noindent\textit{#1}\vspace{2pt}\\
}

\begin{document}

\begin{center}
{\LARGE \textbf{Ben Athiwaratkun}}\\[6pt]
\href{https://linkedin.com/in/benathi}{linkedin.com/in/benathi} \quad
\href{https://benathi.github.io}{benathi.github.io} \quad
\href{mailto:ben.athiwaratkun@gmail.com}{ben.athiwaratkun@gmail.com}
\end{center}

%--------------------------------------------------------------
\section{Summary}
Research and technical leader specializing in frontier-model inference efficiency: speculative decoding, quantization, model architecture, routing, RL rollout efficiency, and kernel generation. Currently Senior Director of Research at Together AI, where my work bridges research and production by inventing algorithms and shipping them into high-throughput inference systems.

%I focus on making LLM inference radically faster and cheaper---treating inference-time compute, not just training, as a primary driver of model capability. In 2023, I saw the growing importance of inference efficiency---while most efforts were still focused on training---and joined Together AI to build an end-to-end inference optimization R\&D system spanning speculative decoding, KV cache compression, and multi-model routing, with research continuously shipping to production.

%--------------------------------------------------------------
\section{Current Role}

\entrytitle{Senior Director of Research (Core ML), Together AI}{2024--present}
\entrysubtitle{Built and lead a 15-person Core ML team focused on frontier-model inference efficiency. Remain hands-on in high-leverage technical work: personally implementing key algorithms in the production inference engine, guiding major design decisions, and co-designing core architectural and algorithmic innovations.}


% for this direction I want to signal that we need to 
\noindent\textbf{Speculative Decoding}
\begin{itemize}
\item
Led the team to develop an in-house algorithm and directly implemented it in the production inference engine, achieving state-of-the-art inference on NVIDIA Blackwell with \href{https://www.together.ai/blog/fastest-inference-for-the-top-open-source-models}{top rankings} across leading open-source models in independent benchmarks. %, following \href{https://www.together.ai/blog/fastest-inference-for-deepseek-r1-0528-with-nvidia-hgx-b200}{earlier results on DeepSeek R1}.
\item
Led the team to investigate various approaches to speculative decoding, including the architecture speed versus acceptance length trade off. Invented model stacking, an architectural adaptation that parallelizes layers to accelerate deep draft models by 80\%, yielding 30\% end-to-end speedup when deployed as a speculator.
\item
Set technical direction for speculative decoding; led the team to develop training and architecture recipes and to ship adaptive speculator systems that learn from live traffic with zero-downtime updates across multiple open-source models.
% \item
% Set research direction on diffusion speculators---framing the kernel-efficiency vs.\ engine-constraint tradeoff and the production-readiness path.

\end{itemize}



\vspace{6pt}
\noindent\textbf{Quantization \& Compression}
\begin{itemize}
\item
Scoped and staffed Together AI's quantization workstream. Directed the team to build ToreQuant, a high-performance quantization engine and recipe library; shipped NVFP4/FP8 and INT4 quantization pipelines with per-module precision control and speculator/MTP support across 15+ frontier dense and MoE checkpoints (DeepSeek, Kimi, GLM, Qwen, MiniMax), powering Together AI's serverless models and customer deployments including Cursor's Composer 2.
\item
Co-developed quantization strategies in close collaboration with inference-engine implementation, aligning precision and scaling choices with kernel and serving-system constraints so research methods translate to production speedups with minimal quality loss.
\item
Driving research on extending quantization reliably to 2-bit weight and KV, targeting next-generation hardware.
\end{itemize}

\vspace{6pt}
\noindent\textbf{Model Routing for Agentic Inference}
\begin{itemize}
\item
Scoped a research direction on efficient routing, leading to an open-source framework, \href{https://arxiv.org/abs/2602.13692}{ThunderAgent}, that performs program-aware KV cache scheduling to increase throughput of agentic workloads.
\item
Led the team to integrate ThunderAgent into the inference system at the router layer (SMG, Dynamo, and equivalent internal framework).
\end{itemize}


%\vspace{6pt}
%\noindent\textbf{Architectures for Efficient Inference}
%\begin{itemize}
%\item
%Co-designed Ladder Residual: residual connections that overlap communication and computation, yielding 29\% tensor-parallel speedup on 70B models (ICML 2025).
%\item
%Drove research adapting autoregressive to diffusion language models; team built the first diffusion LM matching autoregressive quality at 3.1$\times$ throughput.
%\end{itemize}



\vspace{6pt}
\noindent\textbf{RL Rollout Efficiency \& Training-Serving Co-Design}
\begin{itemize}
\item
Directed RL rollout development, including efficient weight transfer via encoded deltas, router replay, and LoRA-MoE; contributed to the open-source training engine (XORL).
\item
Led research-to-production with tech leads on key efficiency bottlenecks: shared-prefix parallel decoding for massively parallel rollouts, \href{https://arxiv.org/abs/2511.13841}{distribution-aware speculative decoding} for long-tail trajectories (up to 50\% rollout time reduction), and \href{https://arxiv.org/abs/2602.06932}{adaptive drafters} that co-evolve with shifting policy distributions.
%\item
%Co-authored training recipes (DeepSWE, DeepCoder) with academic partners; team released the open-source CoderForge dataset.
\end{itemize}





\vspace{6pt}
\noindent\textbf{Kernel Agent Harness}
\begin{itemize}
\item Chartered an early-stage direction on agentic kernel generation for production inference: a multi-agent system that generates, tunes, and curates kernels keyed by operation, shape, and graph-level serving constraints. Scoped first production targets against closed-source dependencies, shape variation, and broader computation-graph constraints.
%\item
%Chartered an early-stage research direction on agentic kernel generation, scoping and staffing the effort against competing priorities. Designed the system as a multi-agent LLM swarm that generates, tunes, and curates a library of kernels keyed by operation, shape, and surrounding-system constraints, drawn from at compile time.
%\item
%Motivated by a gap in standard kernel-gen work, which targets single ops in isolation: this system is built for production realities — closed-source dependencies, shape variation, and constraints from the broader computation graph. Currently pre-deployment, with first production targets being scoped.
\end{itemize}




%--------------------------------------------------------------
\section{Highlighted Publications}

\subsection*{Speculative Decoding}
\noindent\textbf{When RL Meets Adaptive Speculative Training: A Unified Training-Serving System}\\
\textit{ICML 2026} \href{https://arxiv.org/abs/2602.06932}{[pdf]}\\[2pt]
Unified training-serving system that co-evolves speculative drafters with the RL policy, eliminating the train/serve distribution gap that degrades spec dec under shifting policy distributions.

\vspace{6pt}
\noindent\textbf{Beat the Long Tail: Distribution-Aware Speculative Decoding for RL Training}\\
\textit{MLSys 2026} \href{https://arxiv.org/abs/2511.13841}{[pdf]}\\[2pt]
Accelerates RL rollouts by targeting the long-tail of trajectory lengths that dominate wall-clock time. Allocates more aggressive speculative decoding budgets to slow trajectories. Up to 50\% rollout time reduction.

\vspace{6pt}
\noindent\textbf{ATLAS: Adaptive Speculators for Production Inference}\\
\textit{Together AI blog, 2025} \href{https://www.together.ai/blog/adaptive-learning-speculator-system-atlas}{[blog]}\\[2pt]
Production speculative decoding drafters that learn from live traffic with zero-downtime updates, keeping draft models aligned with shifting data distributions across multiple open-source models.

\vspace{6pt}
\noindent\textbf{World's Fastest DeepSeek-R1 Inference on NVIDIA Blackwell}\\
\textit{Together AI blog, 2025} \href{https://www.together.ai/blog/fastest-inference-for-deepseek-r1-0528-with-nvidia-hgx-b200}{[blog]}\\[2pt]
Independently benchmarked top-ranked inference for DeepSeek-R1-0528 on NVIDIA HGX B200, powered by in-house speculative decoding algorithms shipped directly into the production engine.

\subsection*{Quantization \& Compression}
\noindent\textbf{CoQuant: Covariance-Aware Rotation for 2-bit KV Cache Quantization}\\
\textit{Under review at NeurIPS 2026}\\[2pt]
Rotates the KV cache along covariance-informed axes before 2-bit quantization, decorrelating outliers so aggressive quantization preserves quality. Complements Kitty in the 2-bit KV serving roadmap.

\vspace{6pt}
\noindent\textbf{System-Aware 4-Bit KV-Cache Quantization for Real-World LLM Serving}\\
\textit{Under review at COLM 2026}\\[2pt]
Identifies which KV cache quantization methods actually work under real serving constraints (paged memory, fused attention, concurrency). Improves throughput via higher batch size with near-zero accuracy loss.

\vspace{6pt}
\noindent\textbf{Kitty: Accurate and Efficient 2-bit KV Cache Quantization with Dynamic Channel-wise Precision Boost}\\
\textit{MLSys 2026} \href{https://arxiv.org/abs/2511.18643}{[pdf]}\\[2pt]
Pushes KV cache quantization to 2-bit with near-lossless quality by dynamically allocating higher precision to outlier channels. Direct path to next-generation low-bit memory serving.

\vspace{6pt}
\noindent\textbf{Search Your NVFP4 Scales!}\\
\textit{MLSys 2026}\\[2pt]
Smart searching of quantization scales to close the quality gap between NVFP4 and FP8/FP16. Directly applicable to frontier model weight quantization and FP4 attention.

\subsection*{Routing \& Multi-Model Orchestration}
\noindent\textbf{Squeeze Evolve: Unified Multi-Model Orchestration for Verifier-Free Evolution}\\
\textit{Preprint, 2026} \href{https://arxiv.org/abs/2604.07725}{[pdf]} \href{https://squeeze-evolve.github.io/}{[project]}\\[2pt]
Multi-model orchestration for evolutionary inference: allocates strong vs. cheap models by marginal utility at each stage. 3x cost reduction while matching verifier-based methods.

\vspace{6pt}
\noindent\textbf{ThunderAgent: A Simple, Fast and Program-Aware Agentic Inference System}\\
\textit{ICML 2026} \href{https://arxiv.org/abs/2602.13692}{[pdf]}\\[2pt]
Program-aware scheduling for agentic inference. Co-schedules LLM calls, KV caches, and tool execution to increase cache hit rate and throughput across serving and RL workloads.

\subsection*{Architecture \& Inference Algorithms}
\noindent\textbf{Introspective Diffusion Language Models}\\
\textit{Under review at COLM 2026} \href{https://arxiv.org/abs/2604.11035}{[pdf]}\\[2pt]
A new generation paradigm that bridges autoregressive and diffusion models. Uses introspective strided decoding to perform verification and generation in the same forward pass, achieving 3.1x throughput while matching AR quality.

\vspace{6pt}
\noindent\textbf{Opportunistic Expert Activation: Batch-Aware Expert Routing for Faster Decode Without Retraining}\\
\textit{ICML 2026} \href{https://arxiv.org/abs/2511.02237}{[pdf]}\\[2pt]
Batch-aware expert routing for MoE models that accelerates decode latency without retraining. Exploits routing structure across concurrent requests in production serving.

\vspace{6pt}
\noindent\textbf{Ladder-Residual}\\
\textit{ICML 2025} \href{https://arxiv.org/abs/2501.06589}{[pdf]}\\[2pt]
Redesigns residual connections to overlap communication and computation in tensor parallelism. 29\% inference speedup for 70B models across 8 GPUs.

\vspace{6pt}
\noindent\textbf{Bifurcated Attention: Accelerating Massively Parallel Decoding with Shared Prefixes in LLMs}\\
\textit{ICML 2024} \href{https://arxiv.org/abs/2403.08845}{[pdf]}\\[2pt]
Splits attention into shared-prefix and per-sequence GEMMs to eliminate redundant memory IO. Crucial for massively parallel sampling and RL rollouts where thousands of candidate answers are sampled for verification and reward signaling.

\vspace{6pt}
\noindent\textit{Full publication list: \href{https://scholar.google.com/citations?user=KZpZTTQAAAAJ}{Google Scholar} | \href{https://benathi.github.io}{benathi.github.io}}

%--------------------------------------------------------------
\section{Past Experience}

\entrytitle{ML Scientist, AWS AI Labs}{Aug 2019--Dec 2023}
\begin{itemize}
\item Developed bifurcated attention for memory-IO efficient parallel decoding, leading to a patent, production deployment, and an ICML 2024 paper.
\item Built a multi-lingual code generation evaluation framework used widely within the team for benchmarking. ICLR 2023 (Spotlight), patent filed.
\item Invented token alignment for subword completion and shipped it to production in AWS Q. ACL 2024, patent filed.
\item Early advocate for using language models' generalization abilities for few-shot structured prediction, before GPT-3. ICLR 2021 (Spotlight), EMNLP 2022 (Oral).
\item Full-stack contributions to internal language models for code: data processing, tokenizer design, pre-training, inference, and evaluation.
\end{itemize}

%--------------------------------------------------------------
\section{Education}

\entrytitle{PhD in Statistics, Cornell University}{2014--2019}
Focus on machine learning. Advisor: Andrew Gordon Wilson. Research on probabilistic word embeddings and density representations. Interned at AWS and Microsoft Research.

\vspace{4pt}
\entrytitle{BA in Mathematics and Economics, Williams College}{2008--2012}
Magna Cum Laude, Phi Beta Kappa, Sigma Xi. International Physics Olympiad Medalist.

%--------------------------------------------------------------
%\section{Awards}

%Silver Medal, International Physics Olympiad (2007).
%\quad Gold Medal, Thailand Physics Olympiad (2006).
%\quad Top-500, Putnam Mathematical Competition (2012).


\end{document}


%% anecdotes on cool research items we should mention and work it in
%% 


